\title{Byte Latent Transformer: Patches Scale Better Than Tokens}

\begin{abstract}
We present the Byte Latent Transformer ({\textsc BLT}), an innovative byte-level LLM architecture that achieves, for the first time, comparable performance to traditional tokenization-based LLMs at scale, while also offering notable enhancements in inference speed and robustness. 
{\textsc BLT} transforms byte sequences into variable-length patches that act as the fundamental computational units of the model. Patch segmentation leverages next-byte entropy: higher-entropy regions receive greater computational focus and model capacity to address their increased complexity. We conduct the first scaling analysis of byte-level architectures under a fixed computational budget, scaling up to 8B parameters and training on 4T bytes. Our findings affirm that scaling with raw byte inputs—bypassing static vocabularies—is feasible. Training and inference processes benefit from adaptively using larger patches when data is more regular, which further yields qualitative gains in reasoning capabilities and generalization to rare cases. In summary, for a set inference cost, {\textsc BLT} outperforms token-based models in scaling by simultaneously increasing both patch length and model capacity.
\end{abstract}

\section{Introduction}
\label{section:intro}

We present the Byte Latent Transformer~(\textbf{BLT}), an architecture that dispenses with the need for tokenization by processing raw byte sequences directly. BLT is the first model to achieve parity in performance with conventional models that rely on tokenization, while also delivering notable gains in computational efficiency and resistance to errors (\S\ref{sec:robustness}). Current large language models (llms) generally adopt an almost fully end-to-end training paradigm, but rely on a heuristic tokenization phase that maps byte streams to a pre-defined, fixed vocabulary of tokens. This process introduces representational biases, which can manifest as sensitivity to specific domains or modalities~\citep{dagangetting}, increased vulnerability to input perturbations~(\S\ref{sec:robustness}), inadequate modeling of spelling and orthographic properties~\citep{edman2024cute}, and unfairness across different languages~\citep{liang2023xlm,petrov2024language,limisiewicz2024myte}.

Historically, tokenization has been a necessary step, as training llms directly on raw bytes would incur excessive computational expenses at large scales due to extended input sequences~\citep{xue2022byt5}. Previous research has sought to address this by introducing more efficient self-attention mechanisms~\citep{el2020characterbert, clark2022canine} or by developing architectures that do not rely on attention at all~\citep{wang2024mambabyte}~(\S\ref{section:related_work}). Nonetheless, such strategies are mainly advantageous for the training of \textit{smaller models}. When scaling up, the bulk of computation in Transformers arises from the substantial feed-forward network layers applied to every byte, rather than from the attention operations.

To optimize computational resource distribution, we introduce a dynamic, adaptive technique for assembling bytes into \textit{patches}~(\S\ref{section:patching}) and present a novel model structure that integrates both byte-level and patch-level information. In contrast to tokenization, {\textsc BLT} does not employ a predetermined patch vocabulary. Instead, flexible byte groupings are converted into latent representations of patches through streamlined, trainable encoder and decoder modules. Our findings demonstrate that this strategy enables \textit{greater} compute efficiency compared to traditional tokenization-dependent models.

LLMs that rely on tokenization typically assign identical computational resources to every token, regardless of their inherent complexity. This prioritizes uniformity over resource efficiency, given that the compression-driven heuristics underlying token formation do not necessarily align with the true difficulty of prediction tasks. A key innovation in our design is enabling the model to selectively allocate computation to tokens based on necessity. For instance, accurately predicting the terminal characters in most words generally requires less representational capacity—a task that does not benefit from a full-sized transformer—especially when compared to the more intricate prediction involved in selecting the first word of a new sentence. This principle is embodied in the {\textsc BLT} model architecture described in~(\S\ref{section:architecture}), where the model comprises three distinct transformer blocks: two compact local models that operate at the byte level, alongside a larger, global \textit{latent transformer}~(\autoref{fig:arch_high_level}).
In order to inform how bytes should be organized into patches, and thereby guide adaptive computation, {\textsc BLT} segments input data according to the predicted entropy of upcoming bytes, resulting in groupings wherein the information density is relatively homogeneous.

We introduce the first scaling analysis of byte-level models in terms of floating point operation (flop) control, scaling up to 8 billion parameters and processing 4 trillion training bytes. Our results demonstrate that it is feasible to train large-scale models directly on byte sequences, removing the need for fixed-vocabulary tokenization. Overall, {\textsc BLT} achieves comparable training flop-controlled performance\footnote{The computational budget for a model is estimated as the total number of Floating Point OPerations (flops) executed.} to Llama 3, while requiring as much as 50\% fewer flops during inference~(\S\ref{section:scaling}).

Additionally, we provide evidence that operating directly on raw bytes leads to considerable gains in capturing and modeling the rare or idiosyncratic segments of data. {\textsc BLT} models exhibit greater resilience to corrupted or noisy inputs and demonstrate superior abilities at handling character-level phenomena, as shown by improvements in orthographic understanding, phonological processing, and low-resource machine translation tasks~(\S\ref{sec:robustness}).

Moreover, with {\textsc BLT} architectures, it becomes possible to scale both the model size and patch size in parallel, without increasing the inference flop requirements. Using larger patch sizes yields overall computational savings at inference, allowing resources to be shifted towards enlarging the global latent transformer, as this component is invoked less frequently. Through inference flop-controlled scaling studies~(\autoref{fig:fixed_inference_scaling}), we identify much improved scaling properties compared to models based on fixed tokenization schemes.

To conclude, this work offers several key contributions: 1) We present {\textsc BLT}, a byte latent llm framework capable of dynamic compute allocation for enhanced flop efficiency, 2) We demonstrate that our approach attains flop-matched performance comparable to Llama 3 at scales up to 8B parameters, with the flexibility to trade modest drops in evaluation results for flop efficiency improvements reaching 50\%, 3) The {\textsc BLT} architecture enables a novel scalability pathway for llms, facilitating increases in model size under a constant inference budget, and 4) We provide evidence of greater resilience of {\textsc BLT} models to noisy inputs as well as their ability to capture sub-word features in data, an area where token-based llms often fall short. The code for training and inference with {\textsc BLT} is publicly available at~\url{https://github.com/facebookresearch/blt}.

\section{Patching: From Individual Bytes to Groups of Bytes}
\label{section:patching}

Dividing byte sequences into \textit{patches} enables {\textsc BLT} to flexibly assign computational resources according to the surrounding data context. As illustrated in Figure~\ref{fig:patching_types}, a range of techniques exists for grouping bytes into patches. More precisely, a patching function $f_p$ splits an input byte sequence $\pmb{x}=\{x_i,|i=1,\ldots n\}$ of length $n$ into $m < n$ patches, represented as $\pmb{p}=\{p_j|j=1,\ldots,m\}$, by mapping each byte $x_i$ to a binary indicator in \{0,1\}, where 1 designates the start of a new patch. For both token-centric and patch-centric architectures, the primary driver of computational expenditure is the total number of steps undertaken by the core Transformer module; in the case of {\textsc BLT}, this translates to the number of patches formed with the selected patching strategy. As a result, the mean \textit{patch size} becomes the principal determinant of computational overhead in both training and evaluation when employing a given patching approach~(\S\ref{section:flops}).

In what follows, we define three specific patching strategies: segmenting using a fixed byte count per patch~(\S\ref{section:static-patch}), using whitespace characters as patch boundaries~(\S\ref{section:white-patch}), and employing dynamic patching where patch divisions are informed by entropy measures from a compact byte language model~(\S\ref{section:dyn-patch}). We conclude with an exploration of incremental patching and a clarification of how patching fundamentally diverges from standard tokenization~(\S\ref{section:bpe-patch}).

\subsection{Strided Patching Every K Bytes}
\label{section:static-patch}

A common and intuitive method for organizing bytes is to partition them into patches of uniform length $k$, as exemplified by MegaByte~\citep{yu2023megabyte}. Employing a fixed stride allows for both training and inference to be handled with ease, offers a direct approach to modifying the mean patch size, and thereby facilitates the regulation of the computational flop requirements. Nonetheless, this patching approach exhibits notable drawbacks. Primarily, it lacks adaptive compute allocation; that is, computational resources are not directed towards regions of higher necessity. For instance, computational steps may be wasted when only processing whitespace in source code, whereas areas dense in information, such as mathematical notation, may not receive enough compute. In addition, this strategy results in inconsistent, non-contextual segmentation, with identical byte sequences—for example, a word—potentially being split differently each time.

\subsection{Space Patching}
\label{section:white-patch}

\citet{slagle2024spacebyte} introduces a straightforward but powerful enhancement over conventional strided patching: new patches are initiated immediately following any space-like bytes\footnote{
    Space-like bytes include any byte not categorized as a Latin character, numerical digit, or UTF-8 continuation byte. Furthermore, every patch must include at least one byte that does not meet the space-like criterion.}, which frequently correspond to meaningful linguistic boundaries in numerous languages. In Space patching, each word receives the modeling capacity of a latent transformer step (i.e., increased computational effort), thereby standardizing the treatment of words across different sequences and ensuring sufficient flops for more challenging predictions that generally occur after spaces. To illustrate, generating the initial byte for the answer to a query like ``Who composed the Magic Flute?\uline{\hspace{.6em}}'' requires greater model effort than generating the subsequent bytes of the answer, since identifying the first letter filters down the list of plausible answers—making the rest of ``Mozart'' substantially more predictable. Nevertheless, space patching has limitations: it does not generalize smoothly across all languages and data domains, and it inherently lacks the flexibility to alter patch sizes. In the following, we introduce an alternative patching approach inspired by the recognition that initial bytes of words are usually hardest to predict, and which inherently allows dynamic adjustment of patch size.

\subsection{Entropy Patching: Using Next-Byte Entropies from a Small Byte LM}
\label{section:dyn-patch}

Instead of utilizing heuristics based on fixed rules like whitespace, our method adopts a data-driven strategy to detect instances of heightened uncertainty in next-byte predictions. To this end, we present \textit{entropy patching}, wherein patch boundaries are established using entropy estimations.

We train a small byte-level auto-regressive language model on the training data for {\textsc BLT} and compute next byte entropies under the LM distribution $p_{e}$ over the byte vocabulary $\mathcal{V}$:

\begin{equation}
H(x_i) = \sum_{v \in \mathcal{V}} p_{e}(x_i=v|\pmb{x}_{<i}) \log p_{e}(x_i=v|\pmb{x}_{<i})
\end{equation}

We employ two strategies to detect patch boundaries based on the entropies $H(x_i)$. The initial method selects locations where the entropy surpasses a fixed global threshold, as demonstrated in~\autoref{fig:patching}. The alternative method focuses on points exhibiting entropies that significantly exceed the value at the preceding step. This latter technique may also be viewed as pinpointing positions where the typically declining entropy trend within a patch is disrupted.

\begin{align*}
    \text{Global Constraint}\hspace{1cm}H(x_t)& > \theta_g\\
    \text{Approx. Monotonic Constraint}\hspace{1cm}H(x_t)& - H(x_{t-1}) > \theta_r
\end{align*}

Patch boundaries are determined in a streamlined preprocessing phase that takes place as data is being loaded. This approach contrasts with that of~\citet{nawrot-etal-2023-efficient}, who utilize a classifier trained specifically to identify patch boundaries based on entropy measurements. Within our experimental framework~(\S\ref{section:experiments}), we evaluate and contrast these two approaches for separating bytes characterized by low versus high entropy.

\subsection{The Byte-Pair Encoding (BPE) Tokenizer and Incremental Patching}
\label{section:incremental-patching}
\label{section:bpe-patch}

Numerous contemporary llms, such as our reference model Llama 3, employ subword tokenization strategies like bpe~\citep{gage1994new, sennrich-etal-2016-neural}. In this work, we define ``tokens'' as groups of bytes selected from a \textit{finite} vocabulary established before model training, whereas ``patches'' are dynamically aggregated byte sequences that do not rely on a fixed vocabulary. An essential distinction between tokens and patches lies in the model's inability, when using tokens, to interact directly with the original byte-level features.

A significant advancement offered by {\textsc BLT} over tokenization-based approaches lies in its reconceptualization of the balance between vocabulary size and computational requirements. In conventional llms, expanding the vocabulary size results in, on average, larger tokens, which in turn reduces the number of model steps. However, this expansion simultaneously enlarges the final projection layer’s output dimension. Consequently, tokenization-based models face limited flexibility in varying both token length and inference efficiency. As an illustration, Llama 3 raised the average token size from 3.7 to 4.4 bytes, but doing so required a fourfold increase in its embedding table relative to Llama 2.

During generation, {\textsc BLT} must determine at each step whether the current position in the byte sequence marks a patch boundary, as this influences whether additional computation via the Latent Transformer is required. Crucially, this decision must be made without knowledge of any subsequent bytes, since the remainder of the sequence has yet to be generated. Consequently, the patching process cannot rely on future context when determining how to segment the byte sequence. Formally, a patching method $f_p$ is deemed incremental if the following holds:
$$f_p(\pmb{x}_{<i}) = f_p(\pmb{x})_{<i}$$
The bpe approach does not qualify as an incremental patching scheme, because a given prefix may yield different tokenizations depending on the particular sequence that follows, thus violating the incremental property stated above\footnote{Introducing a designated delimiter token for patch boundaries allows bpe to serve as an incremental patching scheme, albeit with an increase in the length of the byte sequence.}.

\section{{\textsc BLT} Architecture}
\label{section:architecture}
The {\textsc BLT} framework consists of a sizable global autoregressive language model designed to process patch-based representations. Accompanying this main model are two lightweight local models: one that transforms sequences of bytes into patch embeddings, and another that reconstructs byte sequences from these patch embeddings~(\autoref{fig:arch_high_level}).

\subsection{Latent Global Transformer Model}

\textit{The Latent Global Transformer} refers to an autoregressive transformer model $\mathcal{G}$ comprised of $l_{\mathcal{G}}$ layers, designed to transform a sequence of latent input patch representations $p_j$ into corresponding output patch representations $o_j$. In this work, the index $j$ indicates individual patches, whereas $i$ is reserved for bytes. The architecture utilizes a block-causal attention mask~\citep{dubey2024llama}, limiting attention such that it extends only up to and including the current patch within a single document. The global model is responsible for the majority of floating-point operations (flops) during both pre-training and inference phases. Consequently, selectively activating this model enables dynamic regulation of computational resources allocated to different segments of the input and output, according to their respective complexities.

\subsection{Local Encoder}
\label{section:local}

The \textit{Local Encoder Model}, symbolized as $\mathcal{E}$, is constructed as a compact transformer model comprising $l_{\mathcal{E}} << l_{\mathcal{G}}$ layers. Its primary function is to swiftly convert a sequence of input bytes $b_i$ into rich patch-level representations $p_j$. Unlike conventional transformer frameworks, this model incorporates a cross-attention mechanism after every transformer layer. This cross-attention operation aggregates information from byte representations into patch-level representations~(\autoref{fig:crossattn}).

Initially, the byte sequence $b_i$ undergoes embedding via a matrix of dimensions $\mathbb{R}^{256 \times h_{\mathcal{E}}}$, producing $x_i$. These vectors can be further enriched through optional hash-embeddings~(\S\ref{sec:hash-emb}). Subsequently, the representations are refined by stacking transformer layers and cross-attention modules in an alternating pattern~(\S\ref{sec:enc-cross-attn}), culminating in the patch representations $p_i$. These patches are forwarded to the global transformer, $\mathcal{G}$, for subsequent processing.

Within the transformer layers, a \textit{local block causal} attention mask is applied such that each byte token considers only a predetermined context window of $w_{\mathcal{E}}$ preceding bytes. This window is permitted to span across moving patch boundaries but remains restricted within a single document. Further specifics about the embedding scheme and the structure of the cross-attention layer are detailed in the following subsections.

\subsubsection{Encoder Hash n-gram Embeddings}
\label{sec:ngram-emb}
\label{sec:hash-emb}
Effectively capturing comprehensive and informative representations at each position $i$ depends heavily on encoding knowledge of the bytes that precede it. Within {\textsc BLT}, this is addressed by representing the byte $b_i$ on its own as well as within the context of a byte n-gram. Specifically, for every position $i$, we initially generate byte-grams

\begin{align}
    g_{i,n}=\{b_{i-n + 1},\ldots, b_i\}
\end{align}

for every byte position $i$ and for all $n$ in the range from three to eight.\footnote{
If $i < n$, we do not consider $n$-grams of size $n$ or greater.}

Next, we employ hash-based $n$-gram embeddings, whereby each byte $n$-gram is assigned, using a hash function, to a specific index in an embedding table $E_{n}^{hash}$ of fixed size—this is done separately for each $n \in\{3,4,5,6,7,8\}$~\citep{bai2010learning}. The computed hash $n$-gram embedding is subsequently combined with the base embedding of the corresponding byte; the sum is then normalized prior to being input into the local encoder. The enriched embedding is defined as

\begin{align}
e_i &= x_i + \sum_{n = 3,...,8}E_{n}^{hash}(\text{Hash}(g_{i,n})) \\
\text{where, Hash}(g_{i,n}) &= \text{RollPolyHash}(g_{i,n}) \% |E_{n}^{hash}|
\end{align}

We standardize $e_i$ by dividing it by the total number of $n$-gram sizes incremented by one, utilizing RollPolyHash as specified in Appendix~\ref{appendix:rollpolyhash}. Section~\ref{section:ablations} presents an analysis of how $n$-gram hash embeddings, across a range of $n$ values and embedding table dimensions, influence trends related to flop-controlled scaling laws. Beyond employing hash-based $n$-gram embeddings, our investigation also encompassed frequency-based $n$-gram embeddings; further methodological specifics regarding this line of inquiry are documented in Appendix~\ref{appendix:ngrams}.

\subsubsection{Encoder Multi-Headed Cross-Attention}
\label{sec:enc-cross-attn}
Our approach largely adopts the input cross-attention module design from the Perceiver architecture~\citep{jaegle2021perceiver}, with a key modification: here, latent variables align with the patch-specific representations rather than a constant pool of latent vectors~(\autoref{fig:crossattn}), and each representation focuses attention exclusively on the bytes forming its associated patch. The architecture uses a query vector for each patch $p_j$, constructed by aggregating the byte-level encodings of that patch through a pooling operation, which is then passed through a linear transformation, $\mathcal{E}_{C} \in \mathbb{R}^{h_{\mathcal{E}} \times (h_{\mathcal{E}} \times U_{\mathcal{E}})}$, where $U_{\mathcal{E}}$ denotes the count of cross-attention heads within the encoder. Letting $f_{\text{bytes}}(p_j)$ represent the byte sequence associated with patch $p_j$, the computation proceeds as follows

\begin{align}
P_{0,j} &= \mathcal{E}_{C}(f_\text{bytes}((p_j)), f~ \text{is a pooling function} \\
P_l &= P_{l-1} + W_o\left(\text{softmax}\left(\frac{QK^T}{\sqrt{d_k}}\right)V\right) \\
\text{where } Q_j &= W_q(P_{l-1,j}), K_i = W_k(h_{l-1,i}), V_i = W_v(h_{l-1,i}) \\
h_l &= \text{Encoder-Transformer-Layer}_l(h_{l-1})
\end{align}

Here, $P \in \mathbb{R}^{n_p \times h_{\mathcal{G}}}$ denotes the $n_p$ distinct patch representations provided to the global model. These representations are constructed by aggregating the byte embeddings $e_i$ associated with each patch $p_j$. The matrices $W_q$, $W_k$, $W_v$, and $W_o$ serve as the learnable projections for generating the queries, keys, values, and output, respectively, with keys and values derived as projections of the byte features $h_i$ from the prior layer (using $e_i$ at the initial layer). The attention mechanism applies a patch-specific masking technique, restricting each query $Q_j$ to consider only those keys and values that belong to bytes in its own patch $j$. Since multi-headed attention is applied to $Q, K$, and $V$ and patch representations generally have a larger size ($h_{\mathcal{G}}$) compared to $h_{\mathcal{E}}$, we represent $P_l$ using several attention heads of dimension $h_{\mathcal{E}}$ during cross-attention, and subsequently concatenate these outputs to obtain the final $h_{\mathcal{G}}$-dimensional representations. Pre-LayerNorm is applied to the queries, keys, and values, and positional embeddings are omitted within this cross-attention layer. The module also incorporates a residual connection encircling the cross-attention component.

\subsection{Local Decoder} 

Analogous to the local encoder, the local decoder $\mathcal{D}$ employs a lightweight transformer architecture, featuring $l_{\mathcal{D}} << l_{\mathcal{G}}$ layers. Its task is to reconstruct the sequence of global patch representations $o_j$ back into raw bytes $y_i$. During this process, the local decoder generates each raw byte conditioned on the sequence of previously decoded bytes. Consequently, it utilizes the hidden states obtained from the local encoder corresponding to the byte-level sequence as inputs. The decoder is structured with $l_{\mathcal{D}}$ layers, which alternate between cross-attention and standard transformer operations. The cross-attention mechanism is applied prior to the transformer sublayer, converting the patch-level information into byte-level features; the subsequent transformer layer then processes the byte representations produced by this cross-attention.

\subsubsection{Decoder Multi-headed Cross-Attention} 

In the decoder cross-attention, the roles of the queries and key/values are interchanged i.e. the byte-representations are now the queries, and the patch representations are now the key/values. The initial byte-representations for the cross-attention are initialized as the byte embeddings from the last encoder layer i.e. $h_{l_{\mathcal{E}}}$. The subsequent byte-representations for layer $l$, $d_{l,i}$ are computed as:

\begin{align}
D_0 &= h_{l_{\mathcal{E}}} \\
B_l &= D_{l-1} + W_o\left(\text{softmax}\left(\frac{QK^T}{\sqrt{d_k}}\right)V\right), \\
  &\text{where } Q_i = W_q(d_{l-1,i}), K_i = W_k(\mathcal{D}_C(o_j)), V_i = W_v(\mathcal{D}_C(o_j)) \\
  D_l &= \text{Decoder-Transformer-layer}_l(B_l)
\end{align}

Here, $W_k$ and $W_v$ are the projection matrices for the key and value, respectively; these operate by applying a linear transformation followed by a splitting operation, $\mathcal{D}_C$, on the final patch representations $o_j$ produced by the global model. The query projection matrix, $W_q$, is applied to the byte-level representations $d_{l-1}$ from the previous layer of the decoder transformer (or to $h_{l_{\mathcal{E}}}$ in the case of the initial layer). $W_o$ denotes the output projection matrix. Consequently, $B \in \mathbb{R}^{h_{\mathcal{D}} \times n_b}$, with $n_b$ indicating the number of output bytes. To obtain the subsequent decoder representations $D_l$, a decoder transformer layer processes the output, $B$, from the cross-attention module. Consistent with the cross-attention mechanism used in the local encoder, we employ multiple attention heads, apply pre-LayerNorm normalization, omit positional embeddings, and include a residual connection surrounding the cross-attention layer.

\section{Experimental Setup}
\label{section:experiments}
We establish meticulously controlled experimental conditions to rigorously compare {\textsc BLT} against models reliant on tokenization, ensuring that {\textsc BLT} does not benefit from potentially longer input sequence contexts.

\subsection{Pre-training Datasets} In this study, all model sizes are pre-trained using two primary datasets: 1) The Llama 2 dataset~\citep{touvron2023llama}, containing 2 trillion tokens sourced from a diverse range of publicly available data, which undergo rigorous cleaning and filtering processes to enhance quality; and 2) {\textsc BLT}-1T, a novel collection consisting of 1 trillion tokens drawn from assorted public datasets, with a portion originating from the pre-training data provided by Datacomp-LM~\citep{li2024datacomp}. The Llama 2 dataset supports scaling law analyses, helping to establish the optimal token count as suggested by~\cite{dubey2024llama} to guide architecture decisions for {\textsc BLT}. In contrast, the {\textsc BLT}-1T dataset is employed for a comprehensive pre-training to enable direct comparison with Llama 3 on downstream benchmarks. Importantly, neither dataset incorporates any data derived from Meta’s products or services. Additionally, baseline tokenizer-based experiments utilize the Llama 3 tokenizer, which has a 128K token vocabulary, offering better baseline results in our trials compared to the Llama 2 tokenizer.

\subsection{Entropy Model} In our experiments, the entropy model is constructed as a byte-level language model that is trained over the same data distribution as the complete {\textsc BLT} model. By default, we utilize a transformer architecture consisting of 100 million parameters, spanning 14 layers, and employing a hidden dimension size of 512, with sliding window attention covering 512 bytes. The other hyperparameter settings are aligned with those used for both our local and global transformers. We also conducted a variety of tests exploring changes in model scale, the size of the receptive field, and alternative architectures, which are detailed in \autoref{section:ablations}. Notably, if the receptive field is sufficiently limited in size, the trained entropy model can be stored as a compact and efficient lookup table.

\subsection{Entropy Threshold and Equalizing Context Length} For models employing entropy-based patching, we determine a patching threshold designed to yield a target average \textit{patch size} over the pretraining dataset. In contrast to tokenization, {\textsc BLT} allows for arbitrary selection of \textit{patch size}, which substantially influences the model's effective context window. To ensure comparable average context lengths—and to prevent models with larger patch sizes from receiving undue benefit—we keep the expected total number of bytes per batch fixed. As a result, when patch size increases, the sequence length per batch is proportionally decreased. On Llama 2 data, we apply an 8k byte context limit, whereas for the {\textsc BLT}-1T corpus we expand this to an average context of 16k bytes, all while consistently using a batch size averaging 16M bytes.

Although the average batch size remains unchanged, dynamic patching techniques result in varying proportions of bytes per patch when assembling data batches. In our approach, the {\textsc BLT} training process aggregates batches by patch count, thereby eliminating the need for padding within the resource-intensive latent transformer stage. This approach guarantees a consistent number of patches across batches. To mitigate sudden increases in memory usage due to exceptionally large patches, byte sequences are either padded or truncated to 12k and 24k bytes for Llama 2 and {\textsc BLT}-1T datasets, respectively, during the training phase.

\subsection{Entropy Model Context}
\label{section:entropy-context}
Our experiments reveal that applying entropy patching tends to produce increasingly larger patches in structured formats, such as multiple choice tasks, which characteristically display high degrees of repetition (refer to patching on an MMLU example in \autoref{fig:mmlu_patching}). These changes arise due to reduced entropy within the recurring segments present in the entropy model context. Consequently, for the extensive evaluation of {\textsc BLT}-Entropy at a patch size of 4.5, we refresh the entropy context by introducing new lines and implement an approximate monotonicity constraint, as this approach mitigates the effects of “entropy drift” caused by context length fluctuations. Notably, this adjustment pertains solely to our entropy computations; our method for selecting the entropy threshold remains unchanged.

\subsection{FLOPs Estimation} 
\label{section:flops}

Our calculation of transformer flops primarily adheres to the methodology presented in Chinchilla~\citep{hoffmann2022training}, which accounts for the feed-forward networks, the qkvo projections within the self-attention mechanism, as well as the evaluation of the attention weights and the output projection. However, a key distinction in our approach is the treatment of the input embedding layer: we posit that it is executed through an efficient lookup operation rather than through dense matrix multiplication, rendering its flop count effectively zero. Consistent with prior research, we assume the backwards pass entails double the flop count relative to the forwards pass.

To compute flops \textit{per byte} for {\textsc BLT} models, we add up the flops for the local encoder transformer, the global latent transformer, and the local decoder transformer, together with the cross attention blocks in the encoder and the decoder:

\begin{align}
    \text{FL}_{\text{{\textsc BLT}}} &= \text{Transf. FL}(h_{\mathcal{G}}, l_{\mathcal{G}}, m=n_{ctx}/n_p,V=0) / n_p\\
   &+\text{Transf. FL}(h_{\mathcal{E}}, l_{\mathcal{E}}, m=w_{\mathcal{E}}, V=0) \\
   &+ \text{Transf. FL}(h_{\mathcal{D}}, l_{\mathcal{D}}, m=w_{\mathcal{D}}, V=256) \\ 
    &+ \text{Cross Attn. FL}(h_{\mathcal{E}}, l_{\mathcal{E}}, m=n_p, r=n_p/k) \cdot (k/n_p) \\ 
   &+ \text{Cross Attn. FL}(h_{\mathcal{D}}, l_{\mathcal{D}}, m=k, r=k/n_p)
\end{align}

Here, $n_{ctx}$ denotes the length of the input sequence measured in bytes, while $n_p$ indicates the patch size. The parameter $r$ represents the proportion of queries relative to key/values. The variable $k$ quantifies the relation between the patch dimension and the byte dimension; specifically, it reflects the number of locally trained model segments that are combined to assemble the full model representation (as illustrated in \autoref{fig:crossattn}, $k=2$). The output projection uses a vocabulary of size $V$, but this is confined to the local decoder. The specific context length used by the attention mechanism, $m$, varies depending on whether it processes a sequence at the byte level or at the patch level. We adjust the flop counts for attention computations in each module to account for these differences. Detailed formulas for computing the Transformer-FLOPs and Cross-Attention FLOPs are found in Appendix~\ref{section:flops-appendix}.

\subsection{Bits-Per-Byte Estimation} Perplexity only makes sense in the context of a fixed tokenizer as it is a measure of the uncertainty for each token. When comparing byte and token-level models, following previous work~\citep{xue2022byt5, yu2023megabyte, wang2024mambabyte}, we instead report Bits-Per-Byte (BPB), a tokenizer independent version of perplexity. Specifically:

\begin{align}
    \text{BPB}(x)&= \frac{\mathcal{L}_{CE}(\pmb{x})}{\ln(2)\cdot n_{\text{bytes}}}
\end{align}

Here, the aggregate cross-entropy loss quantifying the uncertainty associated with $\pmb{x}$ is divided by both the total byte count of $\pmb{x}$ and a normalization factor.

\subsection{Transformer Architecture Hyperparameters} The design of all transformer components in {\textsc BLT}, encompassing both the local and global modules, closely adheres to the Llama~3~\citep{dubey2024llama} specification. Specifically, we employ the SwiGLU activation function~\citep{swiglu} within the feed-forward sublayers, and apply rotary positional embeddings (RoPE)~\citep{RoPe} configured with $\theta=500000$~\citep{xiong2024effective} exclusively to self-attention mechanisms. Layer normalization is performed using RMSNorm \citep{RMSNorm}. For self-attention operations that utilize constant mask patterns such as \textit{block causal} or \textit{fixed-window block causal}, Flash attention~\citep{dao2022flashattention} is leveraged; these are configured with a fixed attention window of 512 tokens where applicable. In contrast, since cross-attention operations utilize dynamically constructed, patch-specific attention masks, we employ Flex Attention\footnote{\url{https://pytorch.org/blog/flexattention}} in order to generate fused computational kernels, which considerably accelerates the training process.

\subsection{{\textsc BLT}-Specific Hyperparameters}

To assess the capabilities of {\textsc BLT} models, we design experiments focusing on two aspects: scaling behavior and downstream task performance. For this purpose, we evaluate model configurations of varying sizes—specifically, 400M, 1B, 2B, 4B, and 8B parameters. Detailed architectural hyperparameters for each model can be found in Appendix~Table~\ref{tab:arch}. 

When initializing the queries for the initial cross-attention layer within the local encoder, max-pooling is employed. All variants of {\textsc BLT} incorporate $500,000$ hashes with one hash function, and utilize n-gram sizes spanning from 3 up to 8. The learning rate is consistently set to $4e-4$ across all models. Our decision to use this learning rate for both token-based and {\textsc BLT} approaches is based on hyperparameter sweeps, specifically within the range $1e-3$ to $1e-4$ for 400M and 1B models, which indicated that this value is optimal. In experiments evaluating scaling on Llama-2 data, we follow batch size recommendations from~\cite{dubey2024llama} or adopt their bytewise equivalents. 

Optimization relies on the AdamW optimizer~\citep{adamw}, where $\beta_1 = 0.9$, $\beta_2 = 0.95$, and $\epsilon=10^{-8}$. A linear warm-up phase of 2000 steps precedes the use of a cosine learning rate decay down to zero. Regularization is provided by a weight decay factor of 0.1, and gradients are globally clipped to a maximum norm of 1.0.

\section{Scaling Trends}
\label{section:scaling}

We provide a comprehensive overview of scaling behaviors observed in byte-level models, with implications for advancing the scaling of {\textsc BLT} models. Our investigation seeks to overcome the shortcomings of earlier studies on byte-level modeling by (a) conducting a comparative analysis of performance trends under compute-optimal training configurations, (b) developing parallel 8B models with substantial training datasets (reaching up to 1T tokens/4T bytes) and assessing their effectiveness on downstream applications, and (c) examining how scaling patterns manifest when inference costs are held constant. Subsequent sections will delve into the distinctive benefits associated with processing byte-sequences directly.

\subsection{Parameter Matched Compute Optimal Scaling Trends}
\label{section:parameter-matched-scaling}
We utilize the Llama 2 dataset to train several \textit{compute-optimal} bpe and {\textsc BLT} models at four parameter scales, spanning from 1B up to 8B. For these configurations, we chart the total training flops against language modeling accuracy on a representative sample drawn from the data mixture. The bpe models are configured with the optimal proportion of model parameters to training tokens, as recommended by Llama~3~\citep{dubey2024llama}. The \textit{compute-optimal} strategy is, in theory, intended to maximize training set performance within a fixed compute limit~\citep{hoffmann2022training}, making it an effective baseline for our evaluations. Correspondingly, each bpe model is paired with a {\textsc BLT} model trained under the same data conditions and employing a Latent Transformer with identical size and architecture as its bpe Transformer counterpart.

As shown in~\autoref{fig:scaling-compute-opt} (right), {\textsc BLT} models consistently equal or exceed the performance of their bpe-based peers, with this pattern persisting as model parameters and computational budgets are increased. To our knowledge, {\textsc BLT} represents the first byte-level Transformer framework to demonstrate comparable scaling behavior to BPE-based architectures under compute-optimal conditions. This finding supports our hypothesis that the optimal parameter-to-compute ratio established for bpe models is similarly applicable to {\textsc BLT}, or is at least closely aligned.

Advancements in architecture alongside dynamic patching are both essential for aligning with bpe scaling trajectories. In ~\autoref{fig:scaling-compute-opt} (left), we present a comparison between space-patching-based architectures and Llama 3. For SpaceByte \citep{slagle2024spacebyte}, we simulate performance using the {\textsc BLT} space-patching variant, omitting both n-gram embeddings and cross-attention mechanisms. While SpaceByte demonstrates gains over Megabyte, it still substantially lags behind Llama 3. Turning to \autoref{fig:scaling-compute-opt} (right), we showcase the cumulative gains achieved by integrating architectural innovations with dynamic patching. When scaled, {\textsc BLT} architectures reach comparable performance to leading tokenizer-based systems like Llama 3.

Additionally, we note that the selection of tokenizer significantly impacts the performance of tokenizer-based models; specifically, models utilizing the Llama-3 tokenizer consistently surpass those employing the Llama-2 tokenizer, even when both are trained on identical datasets.

In summary, the {\textsc BLT} architecture demonstrates performance characteristics that are intermediate between Llama 2 and Llama 3 when employing much larger patch sizes. While the bpe tokenizers of Llama 2 and 3 exhibit mean token lengths of 3.7 and 4.4 bytes, respectively, {\textsc BLT} is capable of maintaining similar scaling behaviors with patch sizes averaging 6 or even 8 bytes. Since inference flop count decreases as average patch size increases, opting for an 8-byte patch size results in approximately 50\% reduction in inference compute requirements. Moreover, models configured with larger patch sizes tend to benefit more as both the data and model sizes are scaled. For example, {\textsc BLT} with an 8-byte patch size initially lags behind bpe Llama 2 at the 1B scale but surpasses it at 7B, indicating that larger patch sizes may yield greater benefits at even larger scales; this also raises the possibility that larger patch sizes could be viable as models and compute budgets continue to expand.

\subsection{Beyond Compute Optimal Task Evaluations}
\label{section:evals}
To gain deeper insight into scaling dynamics, we extend training of an 8B {\textsc BLT} model past the standard compute-optimal regime, utilizing the {\textsc BLT}-1T dataset—an expanded, higher-quality collection—and evaluate the model on a diverse range of established classification and generation tasks. For our assessment, we choose benchmarks that cover common sense reasoning, general world knowledge, and code generation:

\paragraph{Classification tasks} consist of ARC-Easy (0-shot)~\citep{Eval_ARC}, Arc-Challenge (0-shot)~\citep{Eval_ARC}, HellaSwag (0-shot)~\citep{Eval_hellaswag}, PIQA (0-shot)~\citep{Eval_piqa}, and MMLU (5-shot)~\citep{hendrycks2020measuring}. For these benchmarks, we adopt a prompt-based scoring protocol, determining choice likelihoods and computing average accuracy across tasks.

\paragraph{Coding related generation tasks:}  We measure the pass@1 metric on MBPP (3-shot)~\citep{Eval_MBPP} and HumanEval (0-shot)~\citep{Eval_HumanEval}, which serve to test model proficiency in Python code synthesis.

In \autoref{tab:evals}, we present a comparison of three distinct models trained on the {\textsc BLT}-1T dataset: a model based on the Llama 3 tokenizer with byte pair encoding (bpe),\footnote{The Llama 3 tokenizer was selected due to its 128k vocabulary, which demonstrates superior performance relative to the 32k vocabulary in Llama 2.} alongside two versions of the {\textsc BLT} model. One variant applies a space-patching approach ({\textsc BLT}-Space), while the other employs a patching strategy guided by entropy measurements ({\textsc BLT}-Entropy), utilizing an approximate monotonicity constraint and resetting the entropy model’s context at newline characters, as elaborated in~\autoref{section:entropy-context}. All three models are trained within an equivalent compute (flop) budget. For {\textsc BLT}-Entropy specifically, we further adjust the entropy threshold at inference, lowering it from 0.6 to 0.1, a modification that enhances performance on downstream tasks, though at the expense of increased inference steps.

The {\textsc BLT}-Entropy model achieves superior performance compared to the Llama 3 model on 4 out of 7 evaluated tasks, despite both models being trained with an equivalent data volume in bytes. This enhancement can likely be attributed to (1) a more effective utilization of computational resources through the use of dynamic patching, and (2) the explicit modeling at the byte level rather than operating on tokens.

Conversely, {\textsc BLT}-Space lags behind the Llama 3 tokenizer in performance across most tasks, with the exception of one. Nevertheless, due to its larger mean patch size of 6 bytes, it offers a notable decrease in inference flops. For context, both the bpe and entropy-patching based approaches generate an average patch size of about 4.5 bytes on the mixed training dataset. Given an equivalent training computation, the model with the increased patch size processes 30\% more data than the other models, potentially driving {\textsc BLT} further from achieving compute-optimality.

\subsection{Patches Scale Better Than Tokens} {\textsc BLT} models enable parallel scaling of both model and patch sizes without altering the total compute cost during training and inference, nor changing the quantity of training examples. Patch-based architectures uniquely support the unrestricted enlargement of patch size, a property not available to token-based models constrained by their static vocabularies, as outlined in Section~\ref{section:bpe-patch}. Increasing the patch length reduces computational demand, thus permitting additional resources to be directed toward expanding the global latent transformer, which operates at a lower frequency as patch size grows.

A fixed inference scaling investigation is performed to evaluate whether increasing model size, while reducing the number of inference steps on larger patches, yields superior performance compared to smaller models with more inference steps. Utilizing initial models with 400m and 3.6B parameters configured with the Llama 2 tokenizer, we identify models with comparable computational footprints using the Llama 3 tokenizer, as well as {\textsc BLT}-Entropy models employing average patch sizes of 6 and 8 bytes within the training data mixture (refer to~\autoref{tab:fixed-inf-params} for specific model configurations). For models employing a patch size of 8, three encoder layers are utilized as opposed to one. Each of these models is trained across a range of training flop budgets.

\autoref{fig:fixed_inference_scaling} illustrates that {\textsc BLT} models display superior scaling properties relative to tokenization-based models across both inference flop categories. Although BPE-based approaches yield stronger results under limited training budgets, {\textsc BLT} models quickly overtake their performance as the budget increases, with the transition occurring just past the compute-optimal threshold. In practical scenarios, allocating additional resources to pretraining—despite the one-off cost—can result in substantially improved model performance at a constant inference expense. For example, within the 8B model class, exemplified by models such as Llama 3.1, pretraining has been conducted on datasets up to two orders of magnitude larger than what would be considered compute-efficient for the model's scale.

The point at which {\textsc BLT} surpasses token-based models now lies somewhat nearer the compute-optimal threshold in the larger flop category, shifting from 3x down to 2.5x of the optimal compute allocation. The model utilizing an increased patch size of 8 within this larger flop regime also exhibits a sharper improvement trajectory, outperforming other configurations at an earlier stage. As elaborated in Section~\ref{section:parameter-matched-scaling}, performance with larger patch sizes converges towards that of BPE models as overall model scale increases. We partially explain this by the reduction in the proportion of total flops consumed by the byte-level Encoder and Decoder components, whose computational growth lags that of the Latent Transformer as scale increases. For instance, when expanding the model's parameter count by 20x—from 400M to 8B—the quantity of {\textsc BLT} local parameters only approximately doubles. This is a critical consideration since the adoption of larger patch sizes solely increases computational cost within the patch Latent Transformer, leaving the byte-level elements unaffected. Consequently, this accounts for why the {\textsc BLT}-Entropy model with patch size 8 increased from 1.6x to 1.7x the parameter count of Llama 2 when scaling up the model size.

To summarize, our investigation into patch-length scaling reveals that the {\textsc BLT} patch-based model architecture attains superior scaling behavior when both the patch size and the overall model size are scaled up concurrently. These advantageous trends not only continue but also intensify at even greater model scales.

\section{Byte Modeling Improves Robustness} Furthermore, we assess the robustness of the {\textsc BLT} model in contrast to models reliant solely on tokens and therefore lacking explicit byte-level representation, and introduce a method to enable pretrained token-based models to operate at the byte level.
\label{sec:robustness}

\subsection{Character-Level Tasks}

One of the initial incentives for developing byte-level models was their inherent resilience to noise present at the byte level in inputs, as well as their ability to recognize the internal structure of tokens—a capability that tokenizer-based architectures often lack. To assess these attributes, we conduct supplementary experiments on benchmarks designed to evaluate robustness against noisy inputs and character-level comprehension. These benchmarks cover both English and multilingual settings, incorporating digits and phonemes. The outcomes of these experiments are summarized in Table~\ref{tab:char_tasks}.

\paragraph{Noisy Data} To assess the robustness of tokenizer-based models in comparison with {\textsc BLT}, we generate perturbed variants of the benchmark classification tasks outlined in Section~\ref{section:evals}. Five different character-level noise injection techniques are employed to introduce textual variability: (a) \textit{AntSpeak}: Converts the input text into a sequence of uppercase letters separated by spaces. (b) \textit{Drop}: Randomly deletes 10\% of all characters. (c) \textit{RandomCase}: Randomly assigns either uppercase or lowercase to 50\% of the text's characters each. (d) \textit{Repeat}: Randomly repeats 20\% of the characters, with repetitions limited to a maximum of four times per character. (e) \textit{UpperCase}: Changes every character in the text to uppercase. For evaluation, each noise process is applied separately to the prompt, the completion, or both, and the final results are averaged. Table \ref{tab:char_tasks} presents the findings for noised HellaSwag~\citep{Eval_hellaswag}, revealing that {\textsc BLT} consistently demonstrates higher robustness than tokenizer-based models. On average, {\textsc BLT} exceeds the performance of models trained on identical data by 8 points and also surpasses the Llama 3.1 model, which is trained on a substantially larger corpus.

\paragraph{Phonology - Grapheme-to-Phoneme (G2P)} We evaluate {\textsc BLT}'s proficiency in converting sequences of graphemes (individual letters that compose a word) into corresponding phonemic representations (symbols indicating pronunciation). Table \ref{tab:char_tasks} shows performance metrics for the G2P task under a 5-shot scenario, utilizing the Phonology Bench~\citep{suvarna2024phonologybench}. The findings indicate that {\textsc BLT} achieves higher accuracy compared to the Llama 3 1T tokenizer-based model, which serves as our baseline for this evaluation.

\paragraph{CUTE}
In order to measure character-level comprehension, we test {\textsc BLT} on the CUTE benchmark~\citep{edman2024cute}, which encompasses multiple tasks divided into three main groups: compositional understanding, recognition of orthographic similarity, and sequence manipulation skills. This suite is notably difficult for most models that rely on tokenizers; although these models often recognize the orthography of their own tokens, they generally face challenges leveraging this awareness for manipulating text. As displayed in Table~\ref{tab:char_tasks}, {\textsc BLT}-Entropy surpasses both BPE Llama 3 variants by over 25 points on this benchmark. Notably, our system excels at character manipulation tasks, attaining a score of 99.9\% in both spelling-focused evaluations. These considerable gains—despite {\textsc BLT} being trained on merely one-sixteenth the data used by Llama 3.1—suggest that encoding character-level knowledge remains difficult for BPE-based systems. Figure \ref{fig:cute_examples} presents example cases where Llama 3's tokenizer approach falters, whereas {\textsc BLT} succeeds. The only tasks where BPE outperforms are word insertion and deletion. Manipulating entire words may pose more challenges for models operating at the byte level, but the performance disparity is modest, hinting that progressing from character to word manipulations could be more feasible than the reverse. For all tasks, we maintain an identical evaluation protocol, utilizing the original prompts provided by Huggingface. It is possible that BPE models could achieve better results with more sophisticated prompt engineering.

\paragraph{Low Resource Machine Translation} We assess {\textsc BLT} on both directions of translation for twenty-one low-resource languages from six widely spoken language families, each with distinct writing systems, utilizing the FLORES-101 benchmark~\citep{goyal2022flores}. SentencePiece BLEU scores are presented in Table~\ref{tab:flores}. The findings show that {\textsc BLT} yields superior performance over a counterpart trained with the Llama 3 tokenizer, providing an average improvement of 2 BLEU points when translating into English and a 0.5-point gain when translating from English. For common language pairs, the results of {\textsc BLT} are similar to or modestly better than those of Llama 3. Nonetheless, {\textsc BLT} exhibits clear performance gains over Llama 3 across many language pairs belonging to low-resource language families, highlighting the robust generalization of byte-level models when handling rare and diverse byte sequences.

\subsection{Training {\textsc BLT} from Llama 3}
\label{sec:distillation}
In this section, we investigate a methodology that enables {\textsc BLT} models to utilize parameters from already-trained tokenizer-based models, thereby facilitating improved and accelerated convergence during training. Specifically, we initialize the global transformer weights in {\textsc BLT} with those from a pre-trained Llama 3.1 model. Following this initialization, the global transformer parameters are fine-tuned using only one-tenth of the learning rate applied to the local encoder and local decoder components, adopting the optimal training step count determined for Llama 3. Table~\ref{tab:distill} provides a comparative analysis against a standard {\textsc BLT} baseline. The results indicate a substantial performance boost for {\textsc BLT} initialized from Llama 3.1, which exceeds the performance of both the original Llama 3 and baseline {\textsc BLT} models when trained for an equivalent number of flops. Furthermore, in comparison to our {\textsc BLT}-Entropy variant—outlined in Table~\ref{tab:evals}—which underwent training on a much larger dataset (1T tokens), the {\textsc BLT} version derived from Llama 3.1 continues to exhibit superior results on the MMLU benchmark. These findings point to its strong potential for effectively lowering training flop requirements.

This framework enables the transformation of tokenizer-based architectures into tokenizer-free alternatives, thereby repurposing a pre-trained LLaMA 3.1 model as a {\textsc BLT} model. For a thorough evaluation, we list the original LLaMA 3.1, which was trained on 15T tokens, in Table \ref{tab:distill} and benchmark it alongside the {\textsc BLT} derived from LLaMA 3. While the adapted model shows only minor performance reductions on MMLU and HumanEval, other benchmarks reveal more pronounced decreases. These results indicate that additional research is required to fully exploit the pre-trained model’s potential and to enhance its effectiveness, especially by refining the composition of data mixtures and tuning hyperparameters.

\section{Ablations and Discussion}
\label{section:ablations}
Here, we present ablation studies that provide rationale for the architectural decisions underlying {\textsc BLT}, as well as the selection of patching strategy and associated hyper-parameters used for training the {\textsc BLT} 8B parameter model on the {\textsc BLT}-1T dataset.

\paragraph{Entropy Model Hyper-parameters} To evaluate how the size of the entropy model and the length of the context window influence scaling behavior, we conduct training runs with byte-level entropy transformer models whose parameter counts span from 1 million to 100 million and whose context windows range between 64 and 512. The relationship between bits-per-byte (bpb) and training FLOPs is illustrated in the scaling law curves generated from our $400m$ and $1b$ {\textsc BLT} models, both trained on the Llama-2 dataset, as shown in \autoref{fig:entropy_ablation}. Our experiments indicate that increasing either the model size or the context window generally improves scaling; however, these gains become marginal once the number of parameters exceeds 50 million.

\paragraph{Types of Patching}

We perform an ablation study of the four patching schemes described in Section~\ref{section:patching}: (1) Strided Patching with strides of 4 and 6, (2) Whitespace-based patching, (3) BPE Tokenizer patching utilizing the Llama 3 tokenizer, and (4) Entropy-driven patching employing a compact byte-level LLM.

Although dynamic patching shortens the actual sequence length, we ensure that the sequence length is controlled so that the context length remains consistent across all patching approaches. Each model is thus exposed to an equivalent number of bytes per sequence during both training and inference on average, eliminating potential confounders related to context size differences. Figure~\ref{fig:scaling-compute-opt} presents the outcomes of these ablation studies. Notably, all alternative patching methods surpass static patching in performance, with space patching demonstrating competitiveness nearly on par with dynamic entropy based patching.

In \autoref{tab:evals_ablation}, we report benchmark results for the {\textsc BLT} models, providing a comparative analysis between tokenizer-based models, space patching, and entropy-based patching. All models were trained on the Llama 2 dataset with an optimal number of steps as described in~\citep{dubey2024llama}. While space patching is a less complex approach that avoids the computational overhead of dynamically applying an entropy model during training, our findings indicate that the improvements achieved via entropy-based patching in terms of scaling behaviors~(Section~\ref{section:scaling}) are also reflected in downstream evaluation tasks.\footnote{Results for space patching were obtained from earlier experiments that did not include cross-attention, though similar outcome patterns are seen with cross-attention as well.}

\paragraph{Cross-Attention} 
In~\autoref{tab:crossattn_ablations_1b}, we evaluate the impact of integrating cross-attention at different stages within both the encoder and decoder components of {\textsc BLT}. Regarding encoder-side cross-attention, we experiment with initializing the queries using three distinct strategies: (1) assigning an identical learned embedding to all global states, (2) employing a hash embedding derived from the patch bytes, and (3) utilizing a pooled version of the encoder’s hidden representation for the patch bytes at the specific encoder layer.

Our results indicate that integrating cross-attention into the \textit{decoder} yields the highest effectiveness. Incorporating cross-attention within the encoder leads to marginal benefits, which are observable solely when queries are initialized via pooling. Furthermore, cross-attention demonstrates particular utility on Common-Crawl data, with its impact being more pronounced when utilizing larger patch sizes.

\paragraph{n-gram Hash Embeddings} 

We evaluate the effects of varying the n-gram hash embedding vocabulary sizes across settings of 0, 100K, 200K, and 400K, with corresponding outcomes summarized in Table~\ref{tab:ngram_ablations_1b}. Our analysis reveals that hash embeddings contribute positively across every domain examined, with particularly notable gains observed for Wikipedia and Github datasets (showing a 0.04 bpb improvement, in contrast to a 0.01 bpb increase observed after 15k steps at the 8B scale). At 8B, adjusting the hash vocabulary from 500K down to 300K only altered performance by approximately 0.001 bpb over 15k steps, suggesting that hash embeddings are crucial for achieving parity in performance between {\textsc BLT} and traditional tokenizer-based models. Nevertheless, the improvements provided by additional hashes become marginal beyond the 300K threshold. Furthermore, our findings suggest that the benefits from hashes and those from cross-attention largely supplement each other, as each enhances different subsets of the evaluation datasets.

\paragraph{Local Model Hyperparameters} Table \ref{tab:local_layers} presents an ablation of different layer counts within the local encoder and decoder components. With the incorporation of hash n-gram embeddings, {\textsc BLT} achieves strong performance when the encoder is kept minimal—consisting of only a single layer—while the decoder is configured to be comparatively deeper.

\section{Related Work}
\label{section:related_work}

\textbf{Character-Level RNNs:} The domain of Character Language Modeling has long been prominent in neural model research~\citep{sutskever2011generating,mikolov2012subword,graves2013generating}, largely due to its capacity to inherently represent out-of-vocabulary terms without needing external fallback strategies. For instance, \cite{kim2016character} developed an architecture that operates solely on character-level input through a combination of convolutional and highway layers preceding LSTM-based RNNs. Their approach achieved performance on par with the leading RNN language models of the period for English, and even surpassed them in handling morphologically complex languages, underscoring a critical benefit of character-based LLMs. Additionally, \cite{kenter2018byte} leveraged byte-level LSTM architectures for machine comprehension tasks, where these models again outstripped word-level systems on languages such as Turkish and Russian, known for their morphological richness. In a related direction, \cite{zhang2015character} demonstrated that character-level convolutional models can outperform word-level counterparts on certain classification tasks. \citet{chung2019hierarchical} further advanced the field by introducing hierarchical LSTM structures augmented with boundary detectors at each layer, enabling the automatic discovery of latent textual hierarchies and thereby boosting the efficacy of character-level language models. In the context of machine translation, ByteNet from \cite{kalchbrenner2016neural} departed from attention mechanisms by utilizing CNN-based layers directly on character sequences.

\textbf{Character-Level Transformers:} Transformer architectures, leveraging attention mechanisms~\citep{vaswani2017attention} along with subword tokenization techniques~\citep{sennrich-etal-2016-neural}, have markedly advanced the capabilities of neural models in language modeling tasks and standardized evaluations. Nevertheless, the reliance on word or sub-word segmentation instills a predetermined inductive bias about the model's abstraction layer. Seeking to unify the benefits of transformers with the potential exhibited by character-level language modeling, \citet{al2019character} employed highly deep transformer networks enhanced by auxiliary loss functions, resulting in transformer-based models that surpassed preceding LSTM-based character language models. Despite this progress, these models still underperformed when compared to word-level LLMs. Similarly, GPT-2~\citep{radfordgpt2} reported that, particularly for large corpora such as the 1 Billion Word Benchmark, language models utilizing byte-level representations lagged behind those based on word-level tokens in terms of performance.

Although \cite{Choe2019BridgingTG} showed that transformer-based byte-level LLMs can surpass their subword-level counterparts with similar parameter counts, these models incur significantly greater computational costs and require longer training times. In a related effort, \cite{el2020characterbert} introduced CharFormer, a BERT variant that constructs word-level representations using convolutions over character embeddings, achieving performance gains in the medical domain, but also at the expense of heightened computational usage. The work of \cite{clark2022canine} presented CANINE, a 150M-parameter encoder-only model that processes raw character sequences. CANINE integrates a deep transformer stack analogous to the global component of our model, supplemented by a local transformer and strided convolutions to reduce the character input size, resulting in better performance than comparable token-based models (such as mBERT) on various multilingual benchmarks. Byte-level encoder-decoder frameworks were also studied in ByT5~\citep{xue2022byt5}, where the model was not reliant on patching techniques. While ByT5 achieved better resilience to input noise and competitive results compared to models with tokenizers using just one quarter of the training data, the absence of patching necessitated computationally intensive attention over each byte, substantially increasing compute requirements. The challenge in efficiently scaling byte-level models arises from the dramatic increase in input sequence length when operating on bytes directly instead of subwords. More recently, employing the Mamba Architecture~\citep{gu2023mamba}, which is capable of sustaining a constant-sized memory state even across extensive contexts, \cite{wang2024mambabyte} trained a byte-level Mamba model—again forgoing patching—and demonstrated superior performance to byte-level transformer baselines in flop-matched settings with models of approximately 350M parameters, as measured by bits-per-byte across multiple datasets.

\textbf{Patching-based approaches:} Leveraging patching techniques can significantly reduce the high computational costs typically associated with byte-level LLMs, possibly preserving their effectiveness. A number of studies have reported promising outcomes when these methods are applied at a relatively modest scale, both in terms of model parameters and the volume of training data. For instance, \cite{nawrot-etal-2022-hierarchical} explored the impact of static patching through downsampling and upsampling mechanisms, proposing the hourglass transformer architecture. This model demonstrated superior results compared to other byte-level counterparts at the 150M parameter level. Building on these findings, \cite{nawrot-etal-2023-efficient} introduced dynamic patching strategies, incorporating end-to-end trained and tokenizer-supervised boundary predictors, as well as an entropy-driven patching method akin to {\textsc BLT}. Their results indicate that such dynamic techniques can surpass contemporaneous vanilla transformers on language modeling benchmarks with models at 40M parameter size trained on roughly 400M tokens. In a different direction, \cite{lester2024training} examined the benefits of training models on sequences compressed via arithmetic coding—a method that surpasses BPE in terms of achievable compression ratios. Employing an equal-info windows strategy, they managed to outperform byte-level baselines in language modeling scenarios, although their approach fell short compared to subword-level baselines.

Our study builds upon and is most directly influenced by MegaByte~\citep{yu2023megabyte}, a causal decoder-only LLM that performs static, fixed-size patching and concatenation of representations to map bytes into patches. On the decoder side, a local model is employed to convert these patches back to the original byte sequence. Their findings indicate that, at the 1B parameter level and on data comprising approximately 400B bytes, MegaByte achieves parity with models utilizing standard tokenizers. Throughout our experiments, we include an ablation of MegaByte and observe that the use of static patching consistently underperforms compared to state-of-the-art tokenizer-based models that are optimally trained under identical FLOPs constraints. We illustrate that {\textsc BLT} addresses this performance discrepancy. Independently, \cite{slagle2024spacebyte} reach similar conclusions, proposing enhancements to MegaByte such as introducing patching strategies based on whitespaces and other space-character bytes, as well as incorporating a local encoder. They report that, under compute-controlled conditions and within certain domains like Github and arXiv at the 1B parameter scale, their method surpasses tokenized-based transformer models. In our study, we further investigate this architecture, concluding that to truly rival cutting-edge token-based architectures like Llama 3 at the byte level, additional architectural advancements are necessary to support further scaling.

\section{Limitations and Future Work}

In this study, we align our model training steps with the optimal configurations established for Llama~3~\citep{dubey2024llama}, primarily to guide our architectural design decisions. It is important to note, however, that these scaling relationships were originally formulated for BPE-level transformers and may not yield the most efficient (data, parameter size) proportions when applied to {\textsc BLT}. Determining precise scaling laws specific to {\textsc BLT} remains an open problem and a promising direction for subsequent research, as this could reveal even more advantageous scaling behaviors for our model. Furthermore, the majority of our experiments are limited to model sizes up to 1B parameters, so optimal architectural configurations may shift as we extend to 8B parameters or more, potentially unlocking superior performance at larger scales.

Current transformer libraries and codebases are optimized primarily for transformer architectures that operate with tokenizers. Although we provide theoretical experiments where FLOPs are matched and utilize efficient solutions like FlexAttention for layers that differ from the standard transformer design, our implementations may still fall short of tokenizer-based models regarding wall-clock performance, indicating there remains room for additional optimization.

Although {\textsc BLT} employs an independently trained entropy model for the purpose of patching, an intriguing avenue for future research would be to jointly optimize the patching model in an end-to-end manner. Section \ref{sec:distillation} details preliminary studies, which suggest that it is possible to ``byte-ify’’ tokenizer-based models—such as Llama 3, trained on over 10T tokens—by initializing and subsequently freezing the global transformer with their pretrained weights. Continued research along this trajectory could yield approaches that preserve the advantages of bytefying while also advancing performance to exceed that of their tokenizer-based counterparts, potentially obviating the need for retraining from scratch.

\section{Conclusion}
\label{section:conclusion}

This work introduces the Byte Latent Transformer (\textbf{BLT}), an architecture that challenges the traditional reliance on fixed-vocabulary tokenization within large language models. \textsc{BLT} leverages a flexible, learnable mechanism to aggregate bytes into patches, thereby dynamically adapting computation according to the complexity of the input. This approach yields marked enhancements in efficiency and robustness. Through an extensive scaling analysis, we show that \textsc{BLT} can attain performance levels on par with tokenization-dependent models such as Llama 3 for models up to 8B parameters and datasets of up to 4T bytes. Moreover, with only marginal decreases in benchmark scores, \textsc{BLT} can reduce inference flops by as much as 50\%. In addition, \textsc{BLT} opens a new axis for model scaling—enabling simultaneous expansion of both model and patch sizes under fixed inference constraints—a capability particularly well-suited to real-world computational limits. By operating directly on raw bytes, \textsc{BLT} further strengthens robustness against noisy data and enhances the capture of rare or sub-word sequences. Collectively, these findings suggest that \textsc{BLT} offers a viable, efficient, and robust alternative to standard tokenization techniques, laying the groundwork for the next generation of flexible language models.

\end{document}